% Generated draft: requires mathematical and source review.
\documentclass[11pt]{article}
\usepackage[margin=25mm]{geometry}
\usepackage{fontspec}
\setmainfont{Latin Modern Roman}
\usepackage{amsmath,amssymb,mathtools,unicode-math}
\setmathfont{Latin Modern Math}
\usepackage{enumitem,needspace,xurl,hyperref}
\setlength{\parindent}{0pt}
\setlength{\parskip}{5pt}
\setlength{\emergencystretch}{3em}
\newcommand{\N}{\mathbb{N}}
\newcommand{\Z}{\mathbb{Z}}
\newcommand{\Q}{\mathbb{Q}}
\newcommand{\R}{\mathbb{R}}
\newcommand{\C}{\mathbb{C}}
\newcommand{\F}{\mathbb{F}}
\newcommand{\A}{\mathbb{A}}
\newcommand{\PP}{\mathbb{P}}
\newcommand{\E}{\mathbb{E}}
\newcommand{\eps}{\varepsilon}
\newcommand{\dd}{\,\mathrm{d}}
\begin{document}
\section*{UnsolvedMath ID: 30003430}
\textbf{Source problem code:} catalogue entry
\subsection{Definitions and mathematical statement}
For a real-valued portfolio return \(X\) under a probability law \(P\), when \(\mathbb E_P[X^2]<\infty\), its mean and variance are
\[
\mu_P(X)=\mathbb E_P[X],\qquad
\operatorname{Var}_P(X)=\mathbb E_P\!\left[(X-\mu_P(X))^2\right].
\]
These definitions explain the mean--variance terminology; the source does not specify a return variable, investment horizon, probability law, or moment assumptions. Robust portfolio selection considers uncertainty in the inputs or models used to select a portfolio. Here the uncertainty explicitly includes volatility, correlation, and model specification, but the source does not define an uncertainty set or a particular robustness criterion.

\textbf{Source question.} How much can robust mean--variance portfolio strategies improve investor performance under uncertainty in volatility, correlation, and model specification?

A quantitative comparison can be expressed schematically as follows. Once an admissible strategy class \(\mathcal A\), a set \(\Theta\) of possible models, a real-valued performance criterion \(J:\mathcal A\times\Theta\to\R\) oriented so that larger values are preferred, a robust strategy \(a_{\mathrm{rob}}\in\mathcal A\), and a benchmark strategy \(a_{\mathrm{base}}\in\mathcal A\) have been specified, define
\[
\Delta(\theta)=J(a_{\mathrm{rob}},\theta)-J(a_{\mathrm{base}},\theta),\qquad \theta\in\Theta.
\]
The question then concerns the magnitude and circumstances of performance improvement. This display is only a comparison template, not a model or theorem supplied by the source. The record does not specify these objects, whether performance is evaluated under a true model, a worst-case model, or an empirical procedure, or the quantifiers governing any proposed improvement. Consequently, it does not determine a unique mathematical assertion or a universal positivity claim for \(\Delta\).
\subsection{Short English statement}
The problem asks how much robust mean--variance portfolio selection can improve investor performance when volatility, correlations, and the model itself are uncertain. A precise quantitative claim requires an investment model, uncertainty set, portfolio constraints, benchmark, and performance measure. These are absent from the supplied record, so the question remains a research objective requiring further specification.
\subsection{Sources}
\begin{itemize}
\item UnsolvedMath local registry, permanent ID \(30003430\), title: \emph{Performance Gains from Robust Mean--Variance Portfolios}; source problem code \texttt{catalogue entry}. The record belongs to the Oberwolfach Reports Open Problems collection. No individual proposer is identified; the supplied proposed year is 2017.
\item Supplied report citation: \emph{Mathematics of Quantitative Finance (2018)}, Oberwolfach Reports, pp.~683--769, DOI \texttt{10.4171/owr/2017/13}. \url{https://doi.org/10.4171/owr/2017/13}. The background describes Oberwolfach Reports as documenting mathematical workshops and their open-problem discussions. No problem-specific page, speaker, or passage is supplied.
\item Supplied literature reference: Goldfarb--Iyengar, robust portfolio selection. \url{https://doi.org/10.1287/moor.28.1.1.14260}.
\item Supplied literature reference: Robust portfolio optimization comparative analysis. \url{https://pmc.ncbi.nlm.nih.gov/articles/PMC8224264/}.
\item Literature assessment dated \texttt{2026-08-21}: the supplied curation labels the problem \texttt{partially\_solved} and reports that robust mean--variance formulations and distributionally robust extensions are well developed and empirically evaluated, that performance gains depend on uncertainty sets, data, costs, and investor objectives, and that there is no universal improvement theorem. This is supplied dated curation derived from the OWR record, not an independently verified literature assessment; its time-sensitive status and cited claims require independent verification.
\item The registry supplies the statement status \texttt{corrected\_verified}, with an explanation that mathematical content was preserved while notation, typography, and sentence structure were repaired for publication. This is a supplied editorial qualification, not independent verification of the formulation or its mathematical status.
\end{itemize}
\end{document}
